%% kapa: exam=1 round=31 year=2020 period=1 form=A subject=real_estate subject_ko=부동산학원론 seq=13 number=93
%% source: Q-Net 감정평가사 기출문제 (cst00309) artlSeq=5207219 file=extracted/1008/5207219/A형/A형/2020년도 제31회 감정평가사 1차 시험_1교시_A형.pdf
%% conversion: kordoc 4.24.51 -> tools/gen1.py (auto); review=proofread-sonnet; flags=none
\begin{question}{93}
디파스퀠리-위튼(DiPasquale \& Wheaton)의 사분면 모형에 관한 설명으로 옳지 않은 것은? (단, 주어진 조건에 한함)
\begin{choices}
  \item 1사분면에서는 부동산공간시장의 단기공급곡선과 수요곡선에 의해 균형임대료가 결정된다.
  \item 2사분면에서는 부동산의 임대료가 가격으로 환원되는 부동산자산시장의 조건을 나타낸다.
  \item 3사분면에서 신규 부동산의 건설량은 부동산가격과 부동산개발비용의 함수로 결정된다.
  \item 4사분면에서는 신규 부동산의 건설량과 재고의 멸실량이 변화하여야 부동산공간시장의 균형을 이룰 수 있다.
  \item 이 모형은 부동산이 소비재이면서도 투자재라는 특성을 전제로 한다.
\end{choices}
\end{question}
\answer{4}
